\section{Founder 心法：非线性、博弈与时代年龄}

前面讨论产品和 Agent，本章收束到 Founder 心法。肖弘反复强调，世界不是线性外推的。创业者不能只在当前条件下做逻辑推理，而要用博弈方式思考：改变局面、影响他人预期、选择时机、成为重要变量。这个观点贯穿他对资本、平台、模型能力、DeepSeek、Manus 和人生状态的判断。

\begin{figure}[H]
\centering
\includegraphics[width=0.90\textwidth]{founder-game-thinking.png}
\caption{Founder 的博弈思维：不是线性外推，而是让自己成为重要变量。自制概念图，依据 03:09:08--03:15:00 对谈内容整理。}
\end{figure}

\begin{knowledgebox}{读图：逻辑推理和博弈思维的差别}
逻辑推理是在既有条件下找最优解；博弈思维是改变条件本身，让其他参与者重新估计你。Founder 不是旁观预测者，而是局中行动者。
\end{knowledgebox}

\subsection{用时代的年龄思考}

本节解释“时代年龄”。肖弘是 92 年出生，经历移动互联网尾声、2B SaaS、生成式 AI、Agent 爆发。他提醒不要只用生理年龄判断机会，而要用时代的年龄：在某个新技术周期刚开始时，所有人都可能重新站上起跑线。AI 应用创业者需要识别自己处在技术周期的哪个年龄段，是红利期、成熟期还是重构期。

\begin{importantbox}{实践经验：创业者要把自己放进时代坐标}
同样 32 岁，在移动互联网尾声和 Agent 元年完全不同。判断自己是否晚了，不只看年龄，而看新平台、新模型和新用户行为是否刚开始重写规则。
\end{importantbox}

\subsection{为什么不做底层模型}

本节回应一个重要问题。肖弘不做底层模型，不是因为不尊重模型能力，而是因为资源禀赋不同。第一梯队模型公司拓展能力边界，值得尊敬；应用公司更擅长关注用户、承接模型能力、做工作流和品牌。长期看，应用公司可以在特定任务上训练小模型或经济适用模型，但不必从 day one 押注通用底座。

\begin{knowledgebox}{术语消化：模型原厂、应用公司、经济适用模型}
\noindent\begin{tabularx}{\textwidth}{p{0.20\textwidth}p{0.34\textwidth}X}
\toprule
概念 & 含义 & 本期中的作用 \\
\midrule
模型原厂 & 训练通用 foundation model 的公司 & 决定模型能力边界和能力出现时间。 \\
应用公司 & 面向具体用户场景做产品的公司 & 承接模型能力外溢，沉淀场景和品牌。 \\
经济适用模型 & 在特定任务上够用、成本可控的小模型 & 应用公司可逐步补齐局部模型能力。 \\
技术平权 & 通用模型让小团队也能做过去很难的能力 & 独立开发者和小团队获得新机会。 \\
\bottomrule
\end{tabularx}
\end{knowledgebox}

\subsection{Founder 是没得选的}

本节保留最后的创始人状态。肖弘说 Founder 是没得选的：一旦进入局中，就不能像旁观者一样只做理性选择，很多事必须扛住。创业生活会变得昂贵，心理状态、生活成本、团队责任、资本承诺和外部变化都会提高选择的代价。Founder 需要面对贪、嗔、痴，也需要承认自己会被时代、资本和欲望牵引。

\begin{warningbox}{Founder 心法不是鸡汤}
“成为变量”不是鼓励盲目冲动，而是提醒：创业者既要敢改变局面，也要看见自己的欲望、恐惧和误判。博弈思维如果没有自省，会变成赌徒心态。
\end{warningbox}

\subsection{把博弈思维操作化}

本节把抽象心法落到行动。博弈不是玄学，它可以拆成四步：识别局面里的关键参与者，判断他们的动机和约束；找到某个大概率会发生但时间不确定的变化；在可承受成本内提前布置；当事件发生时，用产品、传播和组织动作迅速放大。企业微信 SCRM 是一次这样的博弈，Monica 和 Manus 也是类似逻辑：先判断模型能力和大厂动作，再提前放置产品容器。

\noindent\begin{tabularx}{\textwidth}{p{0.18\textwidth}p{0.34\textwidth}X}
\toprule
步骤 & 具体问题 & EP95 中的例子 \\
\midrule
看玩家 & 谁会改变局面，谁有平台权力 & 腾讯治理外挂，模型原厂释放新能力。 \\
看动机 & 对方为什么迟早会这么做 & 平台要治理生态，模型公司要推进能力。 \\
提前布置 & 我能以多大成本先站位 & 先做企业微信 SCRM，先做 AI 插件和 Manus。 \\
事件发生 & 如何迅速放大正反馈 & 外挂封禁当晚发文，DeepSeek 后加速 Agent。 \\
\bottomrule
\end{tabularx}

\begin{knowledgebox}{课堂提示：博弈思维和用户理解并不冲突}
博弈思维负责判断窗口和位置，用户理解负责把产品做对。只有博弈没有用户，就会变成投机；只有用户没有博弈，就可能错过结构性窗口。
\end{knowledgebox}

\subsection{本章小结}

肖弘的 Founder 心法，是把创业看成动态博弈：外部模型能力在变，大厂判断在变，用户心智在变，资本预期在变。Founder 的任务不是线性预测未来，而是在不稳定系统里建立自己的变量权重。

